\par\Needspace{6\baselineskip}\medskip
\hypertarget{native-112-article083}{}\pdfbookmark[2]{83.}{article-083}
\begin{center}83.\end{center}
Si solutionem eiusdem problematis a computo anguli $\Pi$ inchoare malumus, methodo sequente utemur. Subtrahimus ab aequationum I secunda tertiam, a prima tertiam, a prima secundam, quo pacto tres novas sequentes obtinemus:
\[
\begin{aligned}
&\frac{\dfrac1{r'}-\dfrac1{r''}}{2\sin\tfrac12(N''-N')}=\frac{e}{p}\sin(\tfrac12N'+\mathpunct{\cdot}\tfrac12N''-\Pi)\\[6pt]
\mathrm{II.}\qquad&\frac{\dfrac1r-\dfrac1{r''}}{2\sin\tfrac12(N''-N)}=\frac{e}{p}\sin(\tfrac12N+\tfrac12N''-\Pi)\\[6pt]
&\frac{\dfrac1r-\dfrac1{r'}}{2\sin\tfrac12(N'-N)}=\frac{e}{p}\sin(\tfrac12N+\tfrac12N'-\Pi).
\end{aligned}
\]
Duae quaecunque ex his aequationibus secundum lemma II art. 78 dabunt $\Pi$ et $\dfrac{e}{p}$, unde per quamlibet aequationum I habebuntur etiam $e$ et $p$. Quodsi solutionem tertiam in art. 78, II traditam adoptamus, combinatio aequationis primae cum tertia algorithmum sequentem producit. Determinetur angulus auxiliaris $\zeta$ per aequationem
\[
\operatorname{tang}\zeta=\frac{\dfrac{r'}r-1}{1-\dfrac{r'}{r''}}\cdot\frac{\sin\tfrac12(N''-N')}{\sin\tfrac12(N'-N)}\mathpunct{\raisebox{.45ex}{,}}
\]
\SourceLeaf{113}eritque
\[
\operatorname{tang}(\tfrac14N+\tfrac12N'+\tfrac14N''-\Pi)=\operatorname{tang}(45^0+\zeta)\operatorname{tang}\tfrac14(N''-N).
\]
Permutando locum secundum cum primo vel tertio, duae aliae solutiones huic prorsus analogae prodibunt. Quum hac methodo adhibita formulae pro $\dfrac{e}{p}$ minus expeditae evadant, $e$ et $p$ per methodum art. 80 e duabus aequationum I eruere praestabit. Ceterum ambiguitas in determinatione ipsius $\Pi$ per tangentem anguli $\tfrac14N+\tfrac12N'+\tfrac14N''-\Pi$ ita decidi debebit, ut $e$ fiat quantitas positiva: scilicet manifestum est, pro $e$ valores oppositos prodituros esse, si pro $\Pi$ valores $180^0$ diversi accipiantur. Signum ipsius $p$ autem ab hac ambiguitate non pendet, valorque ipsius $p$ negativus evadere nequit, nisi tria puncta data in parte hyperbolae a Sole aversa iaceant, ad quem casum legibus naturae contrarium hic non respicimus.

Quae ex applicatione methodi primae in art. 78, II post substitutiones operosiores orirentur, in casu praesente commodius sequenti modo obtineri possunt. Multiplicetur aequationum II prima per $\cos\tfrac12(N''-N')$, tertia per $\cos\tfrac12(N'-N)$, subtrahaturque productum posterius a priori. Tunc lemmate I art. 78 rite applicato\footnote{Statuendo scilicet in formula secunda $A=\tfrac12(N''-N')$, $B=\tfrac12N+\tfrac12N''-\Pi$, $C=\tfrac12(N-N'')$.} prodibit aequatio
\[
\begin{aligned}
&\tfrac12\left(\dfrac1{r'}-\dfrac1{r''}\right)\operatorname{cotang}\tfrac12(N''-N')-\tfrac12\left(\dfrac1r-\dfrac1{r'}\right)\operatorname{cotang}\tfrac12(N'-N)\\[4pt]
&\qquad=\frac{e}{p}\sin\tfrac12(N''-N)\cos(\tfrac12N+\tfrac12N''-\Pi).
\end{aligned}
\]
Quam combinando cum aequationum II secunda invenientur $\Pi$ et $\dfrac{e}{p}$, et quidem $\Pi$ per formulam
\[
\begin{aligned}
&\operatorname{tang}(\tfrac12N+\tfrac12N''-\Pi)\\[4pt]
&\qquad=\frac{\dfrac{r'}r-\dfrac{r'}{r''}}{\left(1-\dfrac{r'}{r''}\right)\operatorname{cotang}\tfrac12(N''-N')-\left(\dfrac{r'}r-1\right)\operatorname{cotang}\tfrac12(N'-N)}\mathpunct{\cdot}
\end{aligned}
\]
Etiam hinc duae aliae formulae prorsus analogae derivantur, permutando locum secundum cum primo vel tertio.
